% Figure: partial-fraction decomposition of a rational function.
% r(x) = 1/((x-1)(x+2)) = (1/3)/(x-1) - (1/3)/(x+2). The full curve is
% the sum of two simple-pole terms, each a scaled hyperbola about its
% own pole. Plotted on a window away from the poles so no value diverges.
\begin{figure}[htbp]
\centering
\resizebox{0.86\linewidth}{!}{%
\begin{tikzpicture}[scale=1.0,
  ax/.style={->, draw=engblack, thick},
  tick/.style={draw=engblack!60, thin},
  full/.style={draw=engblue, very thick},
  term/.style={draw=engred, thick, dash pattern=on 3pt off 2pt},
  termb/.style={draw=engamber, thick, dash pattern=on 1pt off 2pt},
  asy/.style={draw=engblack!45, thin, dotted},
  ann/.style={font=\footnotesize, align=center},
]
  % coordinate window: x in [-1.6, 0.6] mapped to cm scale 4x horizontally,
  % y in [-1.2, 1.2] mapped scale 2.4x vertically. Origin at (0,0) cm.
  % We place x-axis at y=0. Use manual transform: X = 4*x, Y = 2.4*y.
  \draw[ax] (4*-1.7, 0) -- (4*0.75, 0) node[below right, ann] {$x$};
  \draw[ax] (0, 2.4*-1.25) -- (0, 2.4*1.3) node[left, ann] {$y$};
  % vertical asymptote lines at x = 1 (outside window) and x = -2 (outside).
  % Within the window (-1.6..0.6) there is no pole, curves stay finite.
  \foreach \x in {-1, 0} {
    \draw[tick] (4*\x, -0.06) -- (4*\x, 0.06);
    \node[ann, below, yshift=-2pt] at (4*\x, 0) {\x};
  }
  \foreach \y in {-1, 1} {
    \draw[tick] (-0.06, 2.4*\y) -- (0.06, 2.4*\y);
    \node[ann, left, xshift=-2pt] at (0, 2.4*\y) {\y};
  }
  % term A: (1/3)/(x-1), finite on the window (x-1 in [-2.6, -0.4])
  \draw[term, domain=-1.6:0.55, samples=80, smooth]
    plot ({4*\x}, {2.4*((1/3)/(\x - 1))});
  % term B: -(1/3)/(x+2), finite on the window (x+2 in [0.4, 2.6])
  \draw[termb, domain=-1.6:0.55, samples=80, smooth]
    plot ({4*\x}, {2.4*(-(1/3)/(\x + 2))});
  % full r(x) = 1/((x-1)(x+2))
  \draw[full, domain=-1.6:0.55, samples=100, smooth]
    plot ({4*\x}, {2.4*(1/((\x - 1)*(\x + 2)))});
  \node[ann, text=engblue, anchor=south east] at (4*0.5, 2.4*0.95)
    {$r(x) = \dfrac{1}{(x-1)(x+2)}$};
  \node[ann, text=engred, anchor=north east] at (4*0.5, 2.4*-0.15)
    {$\dfrac{1/3}{x-1}$};
  \node[ann, text=engamber, anchor=south west] at (4*-1.55, 2.4*0.35)
    {$\dfrac{-1/3}{x+2}$};
\end{tikzpicture}%
}%
\caption[Partial-fraction decomposition]{A rational function written
as the sum of its partial fractions. The function
$r(x) = 1/[(x-1)(x+2)]$ with two simple poles equals
$\tfrac{1/3}{x-1} - \tfrac{1/3}{x+2}$, a sum of two terms, one per
pole, each a scaled hyperbola centred on its own pole. On the window
shown, which sits between the two poles, the blue curve is the exact
sum of the red and amber terms at every point. The decomposition is
the algebra behind the inverse Laplace transform of Chapter 3: each
simple term inverts to one exponential mode, so a transfer function
with two poles produces a response that is the sum of two exponentials,
one decay rate per pole.}
\label{fig:vol02:ch02:partial-fraction}
\end{figure}
